%%%%%%%%%%%%%%%%%%%%%%%%%%%%%%%%%%%%%%%%%%%%%%%%%%%%%%%%%%%%%
\subsubsection{模型建立}

% 下列标题按题型调整。优化模型通常保留；预测、机理、评价模型可替换为
% 学习目标、数据生成假设、状态方程、边界条件或指标体系等。

\paragraph{集合、参数与单位}
设……集合为……，已知参数……表示……，单位为……。

\paragraph{决策变量或状态变量}
定义……。其中，变量的取值范围及现实含义为……。

\paragraph{目标函数}
根据题目对……的要求，建立目标函数
\begin{equation}
    \min/\max\quad Z=\cdots .
\end{equation}
式中各项分别衡量……。采用该目标而非……的原因是……。

\paragraph{约束条件或模型假设}
题目中的……要求对应约束
\begin{equation}
    \cdots .
\end{equation}
该约束保证……。其余约束／边界条件依次为……。

\paragraph{完整模型与性质}
综合目标函数和约束，可将本问写为……模型。该模型具有……性质，
其优势是……，适用边界是……。模型输出中的……直接对应题目要求的……。

%%%%%%%%%%%%%%%%%%%%%%%%%%%%%%%%%%%%%%%%%%%%%%%%%%%%%%%%%%%%%
\subsubsection{模型求解}

\paragraph{算法原理}
由于模型具有……结构，本文采用……算法。其核心思想是……；
算法中的……对应模型变量……，……用于保证约束……。

\paragraph{输入、输出与初始状态}
算法输入为……，输出为……。初始解／初始参数通过……获得。

\paragraph{求解步骤}
% 使用能够复述的具体步骤。复杂算法可改为 algorithm 环境、伪代码或流程图。
\begin{enumerate}[label=\arabic*.,leftmargin=3em]
  \item 根据……构造……；
  \item 计算……并更新……；
  \item 若违反……约束，则通过……恢复可行性；
  \item 根据……接受、拒绝或保留候选；
  \item 当……满足停止条件时，输出……。
\end{enumerate}

\paragraph{参数与停止条件}
关键参数及选取依据见表……。算法在……时停止；为避免将单次启发式结果
误写为全局最优，本文将结果表述为……。

%%%%%%%%%%%%%%%%%%%%%%%%%%%%%%%%%%%%%%%%%%%%%%%%%%%%%%%%%%%%%
\subsubsection{结果与分析}

求解得到表……和图……所示结果。

% 结果段采用“数值—现象—原因—意义”。
由表……可知，……（明确数值）。进一步观察到……（跨数据或跨情景现象）。
该现象可能／可以由……解释；其中……已由……验证，而……仅为合理推测。
这说明模型中的……设计能够……，从而回答了题目关于……的要求。

\paragraph{本问小结}
本问建立了……模型，采用……求解，得到……。该结论属于
严格最优／已验证候选集最优／启发式可行上界／预测或情景结果，其适用范围为……。
